\chapter{Proposición analítica y preguntas de evaluación}
\label{cap:hipotesis}

El estudio distingue una consecuencia algebraica de la regla de intercambio de dos preguntas que requieren simulación. Estas son afirmaciones sobre el micromundo especificado sin comisiones; no son hipótesis sobre mercados observados.

\section{H1: corrección ideal con arbitraje acotado}

Sea $e_t=\log(p_t^D/p_t^C)$ y $D_t=(|e_t|-\gamma)_+$. En una piscina aislada sin comisión, con referencia fija, sin otras órdenes y sin una restricción vinculante de inventario, la orden máxima $q=cx$ cambia el precio logarítmico en $\delta_c=2\log(1+c)$. Como el arbitrajista se detiene en la banda, el exceso de error evoluciona según
\[
  D_{t+1}=\max(0,D_t-\delta_c),\qquad
  N_\gamma=\left\lceil\frac{D_0}{2\log(1+c)}\right\rceil \quad(c>0).
\]
Para $c=0$ no hay corrección finita. Este es un referente mecánico derivado de la ecuación del intercambio de producto constante, no un hallazgo empírico. La simulación sirve para cuantificar las desviaciones respecto de ese referente causadas por otros agentes, el orden de eventos, el horizonte finito y las reglas de inventario especificadas.

\section{H2: retiro, escala de órdenes y respuesta de seguimiento}

La regla E3 existente es un disparador de umbral único, no una política adaptativa estimada: retira el 45\,\% de las reservas en cuanto $|e_t|>0.02$. Los choques ejecutados superan este umbral en todas las corridas E3. Por ello, la pregunta que responden las corridas disponibles es más acotada: ¿cómo cambia esta regla fija de retiro la respuesta simulada con tamaños absolutos de órdenes, y cómo cambia el contraste cuando las órdenes escalan con las reservas o se desactiva la acción de pánico? El desplazamiento máximo original respecto al precio del AMM previo al choque se trata solo como diagnóstico; no mide el sobrepaso frente a la nueva referencia.

\section{H3: valor predictivo de los mapas reducidos}

La comparación reservada pregunta si el mapa ajustado de dos estados mejora la predicción a un paso frente a la persistencia y un AR(1) escalar, y si los términos dependientes de la banda mejoran el error de reproducción recursiva frente a un mapa afín. Esta comparación se refiere a trayectorias generadas por el simulador. La reproducción híbrida recibe la trayectoria futura realizada del CEX, por lo que su resultado no evalúa un pronóstico en tiempo real. No se propone una hipótesis de clasificación de regímenes ni de estabilidad.
